\chapter{Complexity — Unary Languages}
%%% generated by the blueprint pipeline from TCSlib.Complexity.CircuitComplexity.UnaryLanguages
\section{Overview}

Declarations of the Lean module \texttt{TCSlib.Complexity.CircuitComplexity.UnaryLanguages}.

\section{Declarations}

\begin{theorem}[Languages contained in the all-ones language]
\lean{Language.le_allOnes_iff}
\label{Language.le_allOnes_iff}
\leanok
\uses{Language.allOnes, Language.mem_allOnes_iff}
\difficulty{4}
For every language $L$ over the Boolean alphabet, $L$ is contained in the language of
all-ones words if and only if every word in $L$ consists entirely of true bits;
equivalently, each such word is a list of $n$ true bits for some $n$.
\end{theorem}

\begin{definition}[Unary language]
\lean{Language.unary}
\label{Language.unary}
\leanok
\uses{Language.allOnes}
For a set $S$ of natural numbers, the unary language over Boolean words consists of
exactly those words whose entries are all true and whose length belongs to $S$.
\end{definition}

\begin{theorem}[Membership in a unary language]
\lean{Language.mem_unary_iff}
\label{Language.mem_unary_iff}
\leanok
\uses{Language.allOnes, Language.unary}
\difficulty{1}
Let $S$ be a set of natural numbers and let $w$ be a Boolean word. Then $w$ belongs to
the unary language determined by $S$ if and only if all bits of $w$ are true and the
length of $w$ belongs to $S$.
\end{theorem}

\begin{theorem}[Unary languages are contained in the all-ones language]
\lean{Language.unary_le_allOnes}
\label{Language.unary_le_allOnes}
\leanok
\uses{Language.allOnes, Language.unary}
\difficulty{1}
For every set $S$ of natural numbers, every Boolean string whose entries are all true
and whose length belongs to $S$ is a string whose entries are all true; equivalently,
the unary language determined by $S$ is contained in the all-ones language.
\end{theorem}

\begin{theorem}[Membership of unary strings]
\lean{Language.replicate_mem_unary_iff}
\label{Language.replicate_mem_unary_iff}
\leanok
\uses{Language.mem_unary_iff, Language.unary}
\difficulty{3}
For every set $S$ of natural numbers and every natural number $n$, the list of length
$n$ consisting entirely of true entries belongs to the language of all-true Boolean
lists whose lengths lie in $S$ if and only if $n\in S$.
\end{theorem}

\begin{theorem}[Languages encoded by all-ones strings]
\lean{Language.exists_le_allOnes}
\label{Language.exists_le_allOnes}
\leanok
\uses{Language.allOnes, Language.replicate_mem_unary_iff, Language.unary, Language.unary_le_allOnes}
\difficulty{1}
For every set $S$ of natural numbers, there is a language $L$ of Boolean strings
contained in the language of all-ones strings such that, for every natural number $n$,
the length-$n$ all-ones string belongs to $L$ if and only if $n\in S$.
\end{theorem}

\begin{definition}[Boolean NOT gate]
\lean{BoolCircuit.notGateOp}
\label{BoolCircuit.notGateOp}
\leanok
\uses{BoolCircuit.GateOp}
The NOT gate over the two-element Boolean alphabet has one input slot and outputs the
complement of its input, interchanging $0$ and $1$.
\end{definition}

\begin{theorem}[NOT operation belongs to the standard gate set]
\lean{BoolCircuit.notGateOp_mem_stdGateOps}
\label{BoolCircuit.notGateOp_mem_stdGateOps}
\leanok
\uses{BoolCircuit.notGateOp, BoolCircuit.stdGateOps}
\difficulty{1}
The Boolean NOT gate operation on $	ext{Fin }2$, given by the unary function $x\mapsto
1-x(0)$, belongs to the standard set of gate operations.
\end{theorem}

\begin{definition}[Constant-zero circuit nodes]
\lean{BoolCircuit.constZeroNodes}
\label{BoolCircuit.constZeroNodes}
\leanok
For each natural number $n$, define the node type at layer $0$ to be the set of $n$
input positions, and the node types at layers $1$ and $2$ each to be a singleton.
\end{definition}

\begin{definition}[Constant-zero circuit]
\lean{BoolCircuit.constZeroCircuit}
\label{BoolCircuit.constZeroCircuit}
\leanok
\uses{BoolCircuit.LayeredCircuit, BoolCircuit.andGateOp, BoolCircuit.constZeroNodes, BoolCircuit.notGateOp}
For each natural number $n$, there is a depth-two layered circuit from $n$ Boolean
inputs to a singleton output. Its first layer of gates computes a nullary AND, and its
second layer applies NOT, so the output is constantly $0$.
\end{definition}

\begin{theorem}[Constant-zero circuit output]
\lean{BoolCircuit.constZeroCircuit_eval₁}
\leanok
\uses{BoolCircuit.LayeredCircuit.eval₁, BoolCircuit.constZeroCircuit}
\difficulty{1}
For every natural number $n$ and every input assignment $x : \\mathrm{Fin}(n) \\to
\\mathrm{Fin}(2)$, the two-layer circuit with $n$ input nodes and a single output node
evaluates to $0$.
\end{theorem}

\begin{theorem}[Size of the Zero Circuit]
\lean{BoolCircuit.constZeroCircuit_size}
\label{BoolCircuit.constZeroCircuit_size}
\leanok
\uses{BoolCircuit.LayeredCircuit.size, BoolCircuit.constZeroCircuit}
\difficulty{1}
For every natural number $n$, the zero circuit with $n$ inputs has size $2$.
\end{theorem}

\begin{theorem}[Finiteness of the Zero Circuit]
\lean{BoolCircuit.constZeroCircuit_finite}
\label{BoolCircuit.constZeroCircuit_finite}
\leanok
\uses{BoolCircuit.LayeredCircuit.Finite, BoolCircuit.constZeroCircuit}
\difficulty{4}
For every natural number $n$, each layer of the zero circuit with $n$ inputs has a
finite set of nodes.
\end{theorem}

\begin{theorem}[Gate operations in the zero circuit]
\lean{BoolCircuit.constZeroCircuit_onlyUsesGates}
\label{BoolCircuit.constZeroCircuit_onlyUsesGates}
\leanok
\uses{BoolCircuit.LayeredCircuit.onlyUsesGates, BoolCircuit.andGateOp_mem_stdGateOps, BoolCircuit.constZeroCircuit, BoolCircuit.notGateOp_mem_stdGateOps, BoolCircuit.stdGateOps}
\difficulty{4}
For every natural number $n$, the depth-two circuit from $br^n$ to a singleton output
whose first gate layer consists of the empty-input AND gate and whose second gate layer
consists of a NOT gate uses only gate operations from the set comprising the identity,
NOT, and unbounded fan-in AND gates on $\\mathrm{Fin}\,2$.
\end{theorem}

\begin{definition}[Unary language circuit family]
\lean{BoolCircuit.unaryFamily}
\label{BoolCircuit.unaryFamily}
\leanok
\uses{BoolCircuit.LayeredCircuitFamily, BoolCircuit.allOnesCircuit, BoolCircuit.allOnesFamily, BoolCircuit.constZeroCircuit, BoolCircuit.constZeroCircuit_finite}
Given a language $L$ of Boolean words, define a family of layered circuits, one for each
input length $n$. At length $n$, the circuit is the all-ones circuit if the word of $n$
true bits belongs to $L$, and otherwise it is the constant-zero circuit; every layer of
every circuit in the family is finite.
\end{definition}

\begin{theorem}[Unary Family on Accepted Words]
\lean{BoolCircuit.unaryFamily_circuit_of_mem}
\label{BoolCircuit.unaryFamily_circuit_of_mem}
\leanok
\uses{BoolCircuit.allOnesCircuit, BoolCircuit.unaryFamily}
\difficulty{1}
Let $L$ be a language over the Boolean alphabet and let $n\in\mathbb N$. If the word
consisting of $n$ copies of true belongs to $L$, then the circuit of the unary family
for inputs of length $n$ is the all-ones circuit on $n$ inputs.
\end{theorem}

\begin{theorem}[Unary Family on Rejected Words]
\lean{BoolCircuit.unaryFamily_circuit_of_not_mem}
\label{BoolCircuit.unaryFamily_circuit_of_not_mem}
\leanok
\uses{BoolCircuit.constZeroCircuit, BoolCircuit.unaryFamily}
\difficulty{1}
Let $L$ be a language over the Boolean alphabet and let $n\in\mathbb N$. If the word
consisting of $n$ copies of true does not belong to $L$, then the circuit of the unary
family for inputs of length $n$ is the constant-zero circuit on $n$ inputs.
\end{theorem}

\begin{theorem}[Gate Set of the Unary Family]
\lean{BoolCircuit.unaryFamily_onlyUsesGates}
\label{BoolCircuit.unaryFamily_onlyUsesGates}
\leanok
\uses{BoolCircuit.LayeredCircuitFamily.OnlyUsesGates, BoolCircuit.allOnesFamily_onlyUsesGates, BoolCircuit.constZeroCircuit_onlyUsesGates, BoolCircuit.stdGateOps, BoolCircuit.unaryFamily, BoolCircuit.unaryFamily_circuit_of_mem, BoolCircuit.unaryFamily_circuit_of_not_mem}
\difficulty{4}
For any language $L$ of Boolean words, the circuit assigned by the unary family to each
input length uses only identity, NOT, and unbounded fan-in AND gates.
\end{theorem}

\begin{theorem}[Size bound for the unary family]
\lean{BoolCircuit.unaryFamily_size_le}
\label{BoolCircuit.unaryFamily_size_le}
\leanok
\uses{BoolCircuit.LayeredCircuit.size, BoolCircuit.allOnesCircuit_size, BoolCircuit.constZeroCircuit_size, BoolCircuit.unaryFamily, BoolCircuit.unaryFamily_circuit_of_mem, BoolCircuit.unaryFamily_circuit_of_not_mem}
\difficulty{4}
For every language $L$ of Boolean words and every input length $n$, the circuit in the
associated unary family has size at most $2$.
\end{theorem}

\begin{theorem}[Language of the unary family]
\lean{BoolCircuit.unaryFamily_language}
\label{BoolCircuit.unaryFamily_language}
\leanok
\uses{BoolCircuit.LayeredCircuitFamily.language, BoolCircuit.LayeredCircuitFamily.mem_language_iff, BoolCircuit.allOnesFamily_language, BoolCircuit.constZeroCircuit_eval₁, BoolCircuit.unaryFamily, BoolCircuit.unaryFamily_circuit_of_mem, BoolCircuit.unaryFamily_circuit_of_not_mem, Language.allOnes, Language.mem_allOnes_iff}
\difficulty{5}
Let $L$ be a language of Boolean words such that every word in $L$ consists entirely of
true bits. The language recognized by the circuit family that uses the all-ones circuit
at length $n$ exactly when the length-$n$ all-true word belongs to $L$ is equal to $L$.
\end{theorem}

\begin{theorem}[Two-node circuits for sublanguages of all ones]
\lean{Language.inLayeredSIZE_two_of_le_allOnes}
\label{Language.inLayeredSIZE_two_of_le_allOnes}
\leanok
\uses{BoolCircuit.unaryFamily, BoolCircuit.unaryFamily_language, BoolCircuit.unaryFamily_onlyUsesGates, BoolCircuit.unaryFamily_size_le, Language.InLayeredSIZE, Language.allOnes}
\difficulty{1}
Let $L$ be a language of Boolean words such that every word in $L$ consists entirely of
true bits. Then $L$ is recognized by a family of layered circuits using only identity,
NOT, and unbounded fan-in AND gates, with at most two non-input nodes for every input
length.
\end{theorem}

\begin{theorem}[Polynomial-size circuits for sublanguages of all ones]
\lean{Language.inLayeredPPoly_of_le_allOnes}
\label{Language.inLayeredPPoly_of_le_allOnes}
\leanok
\uses{Language.InLayeredPPoly, Language.InLayeredSIZE.inLayeredPPoly, Language.allOnes, Language.inLayeredSIZE_two_of_le_allOnes}
\difficulty{2}
Let $L$ be a language of Boolean words such that every word in $L$ consists entirely of
true bits. Then $L$ is recognized by a family of layered circuits using only identity,
NOT, and unbounded fan-in AND gates, with size bounded by a polynomial in the input
length.
\end{theorem}

\begin{theorem}[Unary languages have polynomial layered size]
\lean{Language.unary_inLayeredPPoly}
\label{Language.unary_inLayeredPPoly}
\leanok
\uses{Language.InLayeredPPoly, Language.inLayeredPPoly_of_le_allOnes, Language.unary, Language.unary_le_allOnes}
\difficulty{1}
For every set $S$ of natural numbers, the language of all-true Boolean strings whose
length belongs to $S$ is computable by a family of finite layered circuits using only
identity, NOT, and unbounded fan-in AND gates, with size bounded by $a(n+1)^k$ for some
natural numbers $a$ and $k$ at every input length $n$.
\end{theorem}

\begin{theorem}[Unary languages have polynomial circuit size]
\lean{Language.unary_inPPoly}
\label{Language.unary_inPPoly}
\leanok
\uses{Language.InLayeredPPoly.inPPoly, Language.InPPoly, Language.unary, Language.unary_inLayeredPPoly}
\difficulty{1}
For every set $S$ of natural numbers, the language of all-true Boolean strings whose
length belongs to $S$ is computable by a family of well-formed Boolean directed acyclic
circuits with fan-in at most two, whose size is bounded by $a(n+1)^k$ for some natural
numbers $a$ and $k$ at every input length $n$.
\end{theorem}
